\section{Introduction}\label{sec:introduction}
An absolute number of coordinate sweeps suffices to mix the full Thorp
permutation. Since a sweep consists of $d$ physical shuffles on $N=2^d$
cards, this gives the optimal order $\Theta(d)$ for the mixing time.
A single sweep already sends each card to a uniform position, but the
joint positions retain dependencies because all cards use the same
random switches. We bound these dependencies through the density of
ordered partial permutations, then use representation multiplicities
to pass from partial information to the full deck.

Write $N=2^d$ with $d\ge1$. The positions are the binary words in $\{0,1\}^d$.
A coordinate layer independently swaps or leaves unchanged the cards on
every edge in one coordinate direction, with probability $1/2$ each.
A sweep visits all $d$ directions in order. In the rotating coordinate
frame this is exactly $d$ physical Thorp shuffles; the precise physical
map is given in Section~\ref{sec:prelim}.
Let $T_N$ be the average action of one sweep. Its adjoint is the average
action of a sweep in the reverse order. Thus $K_N=T_N^*T_N$ is the
positive operator represented by a sweep followed by its reflection,
using fresh independent switches.

For $0\le k\le N$, let $\mathcal X_{N,k}$ be the set of ordered lists of
$k$ distinct positions, let $(N)_k=N!/(N-k)!$, and let $u_{N,k}$ be its
uniform probability measure. For $x,y\in\mathcal X_{N,k}$ define
$R_x(y)=(N)_k\Pr_{K_N}(x\mapsto y)$.
Thus $R_x$ is a density relative to a uniform injection. For a partition
$\lambda\vdash N$, write $D_\lambda$ for its irreducible dimension and
$T_N(\lambda)$ for the sweep's average matrix.

\begin{theorem}[Routing density and representation contraction]\label{fac:main}
There are absolute constants $\eta,c,C>0$ such that, with $a=1/1000$,
\[
 \log\E_{y\sim u_{N,k}}R_x(y)^{1+a}
 \le Ck(k/N)^\eta
 \qquad(1\le k\le N),
\]
uniformly over every starting injection $x$. For every sufficiently
large dyadic $N$ and every nontrivial irreducible representation,
\[
 T_N(\lambda)=0\quad\hbox{or}\quad
 \|T_N(\lambda)\|_{\op}\le D_\lambda^{-c}.
\]
Consequently an absolute number of forward sweeps has worst-start
total-variation distance from uniform tending to zero as $d\to\infty$.
\end{theorem}

The density estimate also measures how much information about the
starting tuple remains. Let $P_x$ be its endpoint law under $K_N$, and
let $H(P_x)$ be its Shannon entropy, with natural logarithms. Jensen's
inequality under $P_x$ gives
\[
 \log (N)_k-H(P_x)=\KL(P_x\|u_{N,k})
 \le a^{-1}\log\E_{u_{N,k}}R_x^{1+a}
 \le (C/a)k(k/N)^\eta.
\]
Thus the entropy deficit per tracked card tends to zero when $k/N\to0$.
At full occupancy the deficit is $O(N)$ after the reflected pair of
sweeps, compared with the initial deficit $\log N!$.

The same saving has a representation-theoretic use. A shape whose first
row is long occurs many times in a moderately larger injection module.
A density cutoff has bounded Hilbert--Schmidt norm on that module, and
these repeated copies force its norm on the shape to be small. Larger
dimensions need a different use of the network, described below. All
traces in the paper are unnormalized; in particular
$\|A\|_{\HS}^2=\Tr(A^*A)$.

Let $q_d$ denote the law of one physical shuffle, let $U_{S_N}$ be the uniform law on $S_N$, and use $\|\mu-\nu\|_{\TV}=\frac12\sum_g|\mu(g)-\nu(g)|$. Define $t_{\mathrm{mix}}(d)$ as the least number of physical shuffles at which this distance is at most $1/4$; by translation the distance is independent of the starting permutation. The lower obstruction counts random bits. Each physical shuffle uses
$N/2$ bits, so after $t$ shuffles at most $2^{tN/2}$ permutations can
occur. The exact bound
\[
 \|q_d^{*t}-U_{S_N}\|_{\TV}\ge1-2^{tN/2}/N!
\]
implies a lower bound $2d-O(1)$ at distance $1/4$.
Together with Theorem~\ref{fac:main}, this gives the optimal order
$\Theta(d)$ in physical time. The theorem does not specify a smallest
mixing constant or a cutoff profile.

\subsection{Network ancestry and previous mixing bounds}
Thorp introduced the shuffle in his study of shuffling and Faro
\cite{Thorp1973}. For power-of-two decks, Morris's 2005 preprint proved
the first polynomial bound, $O(d^{44})$, using evolving sets and a
chameleon process~\cite{Morris2008}. Montenegro and Tetali's
spectral-profile and evolving-set analysis gave $O(d^{29})$
\cite[Theorems 6.14--6.15]{MontenegroTetali2006}. Morris then used
entropy contraction to prove $O((\log N)^4)$ for every even deck size
\cite{Morris2009}, and sharpened the contraction over a sweep to obtain
$O(d^3)$ when $N=2^d$~\cite[Lemma 7 and Theorem 8]{Morris2013}.
These bounds concern the full permutation. Our argument bounds the
joint routing densities and converts them into matrix contraction on
each irreducible representation.

The reflected butterfly has the topology of the classical Bene\v{s}
network~\cite{Benes1964}. Bene\v{s}'s rearrangeability result concerns
which permutations can be routed. Morris already studied the
fair-switch sweep--reflection law under the name \emph{zigzag shuffle}
\cite[Section 4]{Morris2008}. The expanded experiment has $2d$ layers.
Its two central layers average the same subgroup, so their product has
the law of one such average and gives the usual $2d-1$-layer endpoint
description. We retain both layers when counting internal routing
choices.

Gelman and Ta-Shma analyzed this random network through its two-child
recursion~\cite[Section 4, Proposition 1]{GelmanTaShma2014}. Their
Theorem 4 and Lemma 5 give total-variation error at most $k(k-1)/(2N)$
for the image of an unordered $k$-subset. Our estimate concerns ordered
injections and a higher density moment, including dense tuples. The
alternating-color calculation below is the routing structure behind
the recursion; the extra work is to control its cycle counts under
the actual random switch law.

Partial-permutation questions also arise in cryptography. Morris,
Rogaway and Stegers bounded the joint positions of prescribed cards in
their analysis of small-domain encryption
\cite{mrs,MorrisRogawayStegers2018}. Czumaj and V\"ocking proved that,
for every fixed $0<c<1$, the joint positions of any prescribed
$\lfloor cN\rfloor$ cards mix in
$O((\log N)^2)$ physical layers by a non-Markovian coupling
\cite{CzumajVocking2014}. Their full-permutation application pads the
network with empty cards. Czumaj later obtained logarithmic-depth
partial-permutation bounds for other switching networks satisfying an
expansion condition~\cite[Theorems 3.4--3.6]{Czumaj2015}.
The fixed tracked fraction, padding, and choice of network distinguish
these results from mixing the unchanged full-deck Thorp chain.

The passage from Fourier estimates to total variation follows the
finite-group method of Diaconis and Shahshahani
\cite{DiaconisShahshahani1981}. For a conjugacy-invariant law, the
Fourier matrices are scalar and can be estimated through character
ratios. A directed coordinate sweep gives matrix products of
projections, so its singular values and its nonnormal powers must both
be controlled. The inverse-dimension summation used in the final step
is a special case of Liebeck--Shalev
\cite[Proposition 2.5 and Theorem 2.6]{liebeck-shalev2004}; we include an
elementary proof of the estimates needed here. Sharp versions allowing
the exponent to decrease with $N$ are now known
\cite[Proposition 1.8]{teyssier-thevenin2025}. Recent character bounds
also yield cutoff and its profile, relative to the appropriate parity
coset, for nontrivial conjugacy-class walks with a positive fraction of
fixed points
\cite{olesker-taylor-teyssier-thevenin2025,Teyssier2026}.
Those scalar estimates do not provide the matrix contraction required
for a coordinate sweep.

Young branching, hook lengths, and Schur orthogonality connect the
ordered routing densities to individual Fourier blocks
\cite{Sagan2001,Stanley1999,VershikOkounkov2005}.
Later we use the commuting position and label actions on injections
\cite[Section 3.2]{DukesIhringerLindzey2020} and the associated
transposition-content formula for coordinate deletion
\cite[Lemma 3.1 and Theorem 3.5]{AraujoBratten2017}.
These algebraic identities identify the harmonic spaces; the
probabilistic work is to bound the loss of harmonic level when the
network splits.

\subsection{The first proof and the further estimates}
Splitting the reflected network at its outer coordinate leaves two
independent smaller networks. Prescribing the endpoints of some cards
puts two matchings on those cards: one records common input switches,
the other common output switches. Their union consists of alternating
paths and cycles. Each component has two possible assignments to the
children, and each cycle creates one extra factor of two in the route
count. Section~\ref{sec:network} gives the exact density identity and
separates uniform auxiliary colorings from the actual switch law.
This distinction identifies the probability measure under which a
discarded-density tail must be bounded.

Section~\ref{sec:factorial} completes the first proof. We mark several
cycles in distinguished slots, expose their routes, and compare the
number of choices of their starts with the probability that their last
paths close. A contact bound controls the switches already known at
those closures. Factorial moments then bound cycles at each height.
Conditional estimates over alternating height bands sum the bounds
without assuming independence between heights.

The density estimate and the multiplicity cutoff cover shapes with
long first rows and those with very large dimensions. To handle the
remaining dimensions, replace fixed-size central subnetworks by
uniform permutations, allowing identity on a small specified set of
blocks. The modified density moments have arbitrarily small cost per
card. The original positive central operator is bounded by a sum of
these modified operators and an exponentially small remainder.
This completes the dimension-power estimate. Its Fourier consequence
uses submultiplicativity and Schatten norms; a power of $T_N^*T_N$ is
never identified with a power of the generally nonnormal $T_N$.

Section~\ref{sec:certificate} gives an independent certificate argument.
It uses reciprocal cycle lengths to retain explicit density costs for
dense tracked sets, including the margins needed at densities $7/25$
and $1$. Section~\ref{n:strong-providers} instead proves an exponential
cost per tracked card by encoding a permutation while omitting the bits
forced by its closing routes. Its pointwise consequence applies to both
$T_N^*T_N$ and $T_NT_N^*$, for every injection size, after a fixed number
of compositions. Section~\ref{n:coherent-window-route} gives a separate
certificate and entropy-window proof of that density bound.

The first density theorem already implies an $e^{Ck}$ moment bound.
The later density arguments give different mechanisms and lead to
pointwise smoothing and lower-diagram estimates. In particular,
Section~\ref{n:coherent-contact-section} proves a contact-cancellation
bound uniform in the number of tracked cards. Summing out a tuple slot
annihilates a representation at its first injection level; a weighted
forest count turns this cancellation into sparse-level contraction.

The remaining methods retain representation data beyond dimension
alone. Put $k=N-\lambda_1$, $\beta=(\lambda_2,\lambda_3,\ldots)$,
and $h_N(k)=k\log(eN/k)$, with $h_N(0)=0$. The dimension $D_\beta$
detects complexity below the first row that the level $k$ alone cannot
measure. A fixed trace moment controls the combined mass of singular
values, and can therefore be used before the final amplification to
the regular representation. Table~\ref{tab:methods} records these
outputs; its constants are absolute and need not agree between rows.
\begin{table}[ht]
\centering
\small
\begin{tabular}{@{}>{\raggedright\arraybackslash}p{.24\linewidth}>{\raggedright\arraybackslash}p{.66\linewidth}@{}}
\toprule
Estimate & Information retained \\
\midrule
Cycle encoding, Section~\ref{n:strong-providers}
 & $H^u(x,y)\le e^{Ck}/(N)_k$ for either positive orientation;
 $\|T_N(\lambda)\|^2\le e^{Ck}D_\beta^{-a}$. \\
Contact cancellation, Section~\ref{n:coherent-contact-section}
 & $\|T_N(\lambda)\|^2\le\min\{1,(Cdk/N)^{k/2}\}$,
 uniformly for $1\le k<N$. \\
Casimir split, Section~\ref{sec:casimir}
 & $\|K_N(\lambda)\|\le e^{-a h_N(k)}$ and
 $D_\lambda\Tr K_N(\lambda)^{65}\le e^{C h_N(k)}$. \\
Harmonic restriction, Section~\ref{rec:sec-harmonic}
 & $\|T_N(\lambda)\|\le e^{-c\{h_N(k)+\log D_\beta\}}$ and
 $\Tr_{\mathcal X_{N,k}}K_N^{p_0}\le e^{C h_N(k)}$. \\
Adaptive alphabets, Section~\ref{rec:sec-adaptive}
 & $\Tr_{V_\lambda}(T_NT_N^*)^4
 \le e^{-c\log D_\lambda+C h_N(k)}$, paired with level contraction. \\
Core cutoff, Section~\ref{rec:sec-cutoff}
 & Exponentially small overlap with product types of the midpoint
 subcube subgroups whose product dimension is at most
 $D_\lambda^{0.50003}$. \\
\bottomrule
\end{tabular}
\caption{Further density and representation estimates.
All traces are unnormalized. Here $u,p_0$ are absolute integers and
$H$ is either $T_N^*T_N$ or $T_NT_N^*$.}
\label{tab:methods}
\end{table}

The Casimir proof compares the transposition sum with its two child
sums; the sum defect and squared imbalance require different matrix
estimates. The harmonic proof bounds the loss caused by removing
slots, then transfers cycle moments to either lower-diagram dimension
or tuple trace. These arguments share the elementary tensor-swap
estimate, but use different level recursions. The alphabet proof
changes its tensor realization as descendant diagrams shrink and pays
the resulting multiplicity costs across the tree. The core-cutoff
proof chooses between a diagram bound and a bound using the actual
trace at the current recursion frontier. Its cutoff is a projection
onto small product dimensions. Each route supplies its own final
passage to the full permutation law using the common Fourier
preliminaries.

The companion papers use several of these interfaces while retaining
their own arguments. The compatibility-entropy paper uses the
shared-switch contact bound in Lemma~\ref{n:butterfly-contacts} in a
sparse path second-moment estimate
\cite[Lemma 3.1]{CompatibilityEntropy2026}. Prescribed-type signed
tensor moments and one-sided projection-overlap bounds are developed
separately in~\cite[Theorems 1.1 and 3.2]{SignedTensorBudgets2026}.
The row--column geometry paper combines Theorem~\ref{n:strong-density}
and Proposition~\ref{n:strong-tail} with a coherent-overlap argument to
obtain a dimension-weighted fixed moment
\cite[Theorem 10.1]{RowColumnGeometry2026}. The random-subspace paper
applies Corollary~\ref{n:strong-smoothing}, in both positive
orientations, in its Proposition 2.1; complex random-plane tests then
give a regular trace at most $1+N^{-10}$
\cite[Theorem 1.1]{RandomSubspace2026}. An independent
conditional-marginal proof gives the explicit bound of $1600d$
physical shuffles for total-variation distance tending to zero
\cite[Theorem 1.1]{OptimalThorpMixing2026}.
